\documentclass[12pt,a4paper]{article}

% ── Packages ──────────────────────────────────────────────────────────────────
\usepackage[T1]{fontenc}
\usepackage[utf8]{inputenc}
\usepackage{lmodern}
\usepackage[margin=2.5cm]{geometry}
\usepackage{amsmath,amssymb,amsthm}
\usepackage{graphicx}
\usepackage{booktabs}
\usepackage{array}
\usepackage{xcolor}
\usepackage{hyperref}
\usepackage{listings}
\usepackage{caption}
\usepackage{subcaption}
\usepackage{float}
\usepackage{enumitem}
\usepackage{fancyhdr}
\usepackage{titlesec}
\usepackage{mdframed}
\usepackage{multirow}
\usepackage{microtype}
\usepackage{setspace}

% ── Colours ───────────────────────────────────────────────────────────────────
\definecolor{darkblue}{RGB}{31,78,121}
\definecolor{midblue}{RGB}{44,123,182}
\definecolor{lightblue}{RGB}{232,241,250}
\definecolor{darkred}{RGB}{180,30,30}
\definecolor{codegray}{RGB}{245,245,245}
\definecolor{codegreen}{RGB}{0,110,0}
\definecolor{codepurple}{RGB}{100,0,128}

% ── Hyperref ──────────────────────────────────────────────────────────────────
\hypersetup{
  colorlinks=true, linkcolor=darkblue, urlcolor=midblue, citecolor=darkblue,
  pdftitle={Credit Scoring and Portfolio Loss Analysis}
}

% ── Section formatting ────────────────────────────────────────────────────────
\titleformat{\section}{\Large\bfseries\color{darkblue}}{\thesection}{1em}{}
  [{\color{midblue}\titlerule[1.5pt]}]
\titleformat{\subsection}{\large\bfseries\color{darkblue}}{\thesubsection}{1em}{}
\titleformat{\subsubsection}{\normalsize\bfseries\color{darkblue}}{\thesubsubsection}{1em}{}
\titlespacing{\section}{0pt}{18pt}{8pt}
\titlespacing{\subsection}{0pt}{12pt}{4pt}

% ── Header / Footer ───────────────────────────────────────────────────────────
\pagestyle{fancy}\fancyhf{}
\renewcommand{\headrulewidth}{0.4pt}
\renewcommand{\headrule}{\hbox to\headwidth{\color{midblue}\leaders\hrule height \headrulewidth\hfill}}
\setlength{\headheight}{24pt}
\fancyhead[L]{\small\color{gray}\itshape Credit Scoring \& Portfolio Loss Analysis}
\fancyhead[R]{\small\color{gray}\itshape June 2026}
\fancyfoot[C]{\small\color{gray}\thepage}

% ── Code listings ─────────────────────────────────────────────────────────────
\lstset{
  backgroundcolor=\color{codegray}, basicstyle=\ttfamily\small,
  keywordstyle=\color{darkblue}\bfseries, commentstyle=\color{codegreen}\itshape,
  stringstyle=\color{codepurple}, breaklines=true, frame=single,
  framerule=0.4pt, rulecolor=\color{midblue}, numbers=left,
  numberstyle=\tiny\color{gray}, xleftmargin=1.5em, framexleftmargin=1.5em,
  tabsize=4
}

% ── Environments ──────────────────────────────────────────────────────────────
\newmdtheoremenv[backgroundcolor=lightblue, linecolor=midblue, linewidth=1.5pt,
  innerleftmargin=10pt, innerrightmargin=10pt,
  innertopmargin=6pt, innerbottommargin=6pt]{result}{Key Result}[section]

\newmdenv[backgroundcolor=lightblue, linecolor=midblue, linewidth=1.5pt,
  innerleftmargin=12pt, innerrightmargin=12pt,
  innertopmargin=8pt, innerbottommargin=8pt]{insight}

% ── Caption style ─────────────────────────────────────────────────────────────
\captionsetup{font={small,color=darkblue}, labelfont=bf, labelsep=period, skip=4pt}

% ── Graphics path ─────────────────────────────────────────────────────────────
\graphicspath{{../plots/}}

% =============================================================================
\begin{document}
% =============================================================================

% ── Title Page ────────────────────────────────────────────────────────────────
\begin{titlepage}\centering
  \vspace*{2cm}
  {\color{darkblue}\rule{\textwidth}{2pt}}\\[0.4cm]
  {\Huge\bfseries\color{darkblue} Credit Scoring \&\\[0.3cm]Portfolio Loss Analysis}\\[0.4cm]
  {\color{midblue}\rule{\textwidth}{2pt}}
  \vspace{1.2cm}
  {\Large\itshape Lending Club Loan Data --- Statistical Classification \& Basel II Risk Metrics}
  \vspace{0.8cm}\\
  {\large\color{gray} C++17 ML Engines \quad\textbullet\quad Monte Carlo Portfolio Simulation \quad\textbullet\quad R Statistical Analysis}
  \vspace{2.5cm}

  \begin{tabular}{rl}
    \textbf{Dataset:}      & Lending Club loan data (10,000 loans, 14 features) \\[4pt]
    \textbf{Models:}       & LDA, $k$-NN ($k=15$), Logistic Regression (L2) \\[4pt]
    \textbf{Portfolio:}    & One-factor Gaussian copula, 100,000 MC scenarios \\[4pt]
    \textbf{Risk metrics:} & EL, UL, VaR 95/99\%, CVaR, Basel II Capital \\[4pt]
    \textbf{Tools:}        & C++17 (GCC), R 4.3, ggplot2, ragg \\[4pt]
    \textbf{Date:}         & June 2026
  \end{tabular}

  \vfill{\color{midblue}\rule{0.5\textwidth}{0.8pt}}
\end{titlepage}

\tableofcontents\newpage

% =============================================================================
\section{Introduction}
% =============================================================================

Credit risk is the possibility that a borrower fails to meet their obligations,
causing the lender a financial loss. It is one of the oldest and most studied
problems in quantitative finance, and underlies the Basel II/III regulatory
capital framework that governs how banks provision against loan portfolios.

This project addresses credit risk from two complementary angles:
\begin{enumerate}[leftmargin=2em]
  \item \textbf{Borrower-level scoring}: classify each loan as likely to default
        or not, using discriminant analysis, nearest-neighbour, and logistic
        regression classifiers built from scratch in C++.
  \item \textbf{Portfolio-level risk}: aggregate individual predicted default
        probabilities into a portfolio and simulate the full loss distribution
        via Monte Carlo, yielding VaR, CVaR, and the Basel II capital requirement.
\end{enumerate}

\subsection{Dataset}

The dataset follows the Lending Club loan schema: 10,000 unsecured consumer
loans drawn from the 2007--2011 cohort structure, with 21.2\% historical default
rate. Features include borrower credit grade, FICO score, debt-to-income ratio,
loan amount, interest rate, employment length, and delinquency history. The
target variable is binary: \texttt{Charged Off} (default) vs \texttt{Fully Paid}.

% =============================================================================
\section{Theoretical Background}
% =============================================================================

\subsection{Credit Scoring}

Credit scoring assigns a numerical score or probability to each borrower
reflecting their likelihood of default. Three classifiers are implemented:

\subsubsection{Linear Discriminant Analysis (LDA)}

Fisher's LDA finds the projection direction $\mathbf{w}$ that maximises
between-class separation relative to within-class scatter:
\begin{equation}
  \mathbf{w} = S_W^{-1}(\boldsymbol{\mu}_1 - \boldsymbol{\mu}_0)
  \label{eq:lda}
\end{equation}
where $S_W = \sum_{k\in\{0,1\}}\sum_{i:y_i=k}(\mathbf{x}_i - \boldsymbol{\mu}_k)(\mathbf{x}_i - \boldsymbol{\mu}_k)^\top$
is the pooled within-class scatter matrix. A loan is classified as default if
$\mathbf{w}^\top\mathbf{x} \geq \theta$ where $\theta$ is the midpoint of the
two projected class means.

\subsubsection{$k$-Nearest Neighbours}

The $k$-NN classifier assigns $\hat{P}(\text{default}|\mathbf{x})$ as the
fraction of the $k$ nearest training observations (by Euclidean distance in
standardized feature space) that defaulted:
\begin{equation}
  \hat{p}(\mathbf{x}) = \frac{1}{k}\sum_{i \in \mathcal{N}_k(\mathbf{x})} y_i
  \label{eq:knn}
\end{equation}
We use $k=15$ with an asymmetric decision threshold of 0.35 (rather than 0.50),
reflecting the higher cost of false negatives (missed defaults) in lending.

\subsubsection{Logistic Regression with L2 Regularisation}

The logistic model specifies:
\begin{equation}
  P(Y=1|\mathbf{x}) = \sigma(\mathbf{w}^\top\mathbf{x} + b)
  = \frac{1}{1 + e^{-(\mathbf{w}^\top\mathbf{x} + b)}}
  \label{eq:logistic}
\end{equation}
Parameters are estimated by minimising the regularised cross-entropy loss:
\begin{equation}
  \mathcal{L}(\mathbf{w}, b) = -\frac{1}{n}\sum_{i=1}^{n}
  \left[y_i\log\hat{p}_i + (1-y_i)\log(1-\hat{p}_i)\right]
  + \frac{\lambda}{2}\|\mathbf{w}\|^2
  \label{eq:logloss}
\end{equation}
via mini-batch stochastic gradient descent ($\eta=0.05$, $\lambda=10^{-3}$,
batch size 64, 300 epochs). The L2 penalty prevents coefficient blow-up and
improves generalisation.

\subsection{Evaluation Metrics}

For a binary classifier with threshold $\tau$:
\begin{equation}
  \text{Precision} = \frac{TP}{TP+FP}, \quad
  \text{Recall}    = \frac{TP}{TP+FN}, \quad
  F_1              = 2\cdot\frac{\text{Precision}\cdot\text{Recall}}{\text{Precision}+\text{Recall}}
  \label{eq:metrics}
\end{equation}

The \textbf{AUC-ROC} (area under the receiver operating characteristic curve)
is threshold-independent and measures the probability that the model ranks a
randomly chosen defaulter above a randomly chosen non-defaulter.

\subsection{Portfolio Credit Risk}

\subsubsection{Expected and Unexpected Loss}

For a loan with probability of default $p$, loss given default $L$ (fraction of
exposure), and exposure at default $E$:
\begin{equation}
  \text{EL} = p \cdot L \cdot E, \qquad
  \text{UL} = \sqrt{p(1-p)}\cdot L \cdot E
  \label{eq:el_ul}
\end{equation}
For a portfolio of $n$ correlated loans:
\begin{equation}
  \text{EL}_\text{port} = \sum_{i=1}^n p_i L_i E_i, \qquad
  \text{UL}_\text{port} = \text{std}\!\left[\sum_i L_i E_i \cdot \mathbf{1}[\text{default}_i]\right]
  \label{eq:port_loss}
\end{equation}

\subsubsection{One-Factor Gaussian Copula}

Under the Basel II one-factor model, each borrower $i$ has an asset return:
\begin{equation}
  X_i = \sqrt{\rho_i}\,Z + \sqrt{1-\rho_i}\,\varepsilon_i,
  \quad Z, \varepsilon_i \stackrel{\text{iid}}{\sim} \mathcal{N}(0,1)
  \label{eq:copula}
\end{equation}
where $Z$ is the common systematic factor and $\varepsilon_i$ is idiosyncratic.
Borrower $i$ defaults in a given scenario if $X_i < \Phi^{-1}(p_i)$. The
asset correlation is given by the Basel II formula:
\begin{equation}
  \rho(p) = 0.12\cdot\frac{1-e^{-50p}}{1-e^{-50}} + 0.24\cdot\left[1-\frac{1-e^{-50p}}{1-e^{-50}}\right]
  \label{eq:rho}
\end{equation}

\subsubsection{Value at Risk and Expected Shortfall}

\begin{equation}
  \text{VaR}_\alpha = \inf\{l : P(L > l) \leq 1-\alpha\}, \qquad
  \text{CVaR}_\alpha = \mathbb{E}[L \mid L > \text{VaR}_\alpha]
  \label{eq:var}
\end{equation}
CVaR (also called Expected Shortfall) is a coherent risk measure and is
preferred over VaR for internal capital purposes.

\subsubsection{Basel II Capital Requirement}

The Unexpected Loss that banks must hold capital against:
\begin{equation}
  K(p) = L\left[\Phi\!\left(\frac{\Phi^{-1}(p) + \sqrt{\rho}\,\Phi^{-1}(0.999)}{\sqrt{1-\rho}}\right) - p\right]
  \label{eq:basel}
\end{equation}
where $\Phi^{-1}(0.999) \approx 3.09$ reflects the 99.9\% confidence level of
the Basel internal ratings-based (IRB) approach.

% =============================================================================
\section{Implementation}
% =============================================================================

\subsection{C++ Architecture}

The C++ engine is structured around four headers and a single main driver:

\begin{center}
\renewcommand{\arraystretch}{1.3}
\begin{tabular}{@{}ll@{}}
\toprule
\textbf{File} & \textbf{Responsibility} \\
\midrule
\texttt{include/matrix.h}       & Dense matrix, LU inversion, vector utilities \\
\texttt{include/data\_loader.h} & CSV parsing, feature encoding, standardization, train/test split \\
\texttt{include/classifiers.h}  & LDA, $k$-NN, Logistic Regression, metric computation \\
\texttt{include/portfolio\_mc.h}& One-factor Gaussian copula MC, Basel II formulas, VaR/CVaR \\
\texttt{include/results\_writer.h}& CSV export for all results \\
\texttt{src/main.cpp}           & Pipeline driver: load $\to$ train $\to$ evaluate $\to$ portfolio $\to$ MC \\
\bottomrule
\end{tabular}
\end{center}

\subsubsection{Feature Engineering}

Fourteen numerical features are extracted from the raw CSV. Categorical
variables (grade, home ownership, purpose, term) are label-encoded.
All features are \textbf{standardized} (zero mean, unit variance computed on
the training set and applied to test set) so that LDA's $S_W^{-1}$ is
well-conditioned and logistic regression coefficient magnitudes are comparable.

\subsection{Compilation}

\begin{lstlisting}[language=bash, caption={Compile and run}]
cd src/
g++ -O2 -std=c++17 -I../include main.cpp -o credit_model
./credit_model       # writes all CSVs to ../data/

cd ../R/
Rscript analysis.R   # generates plots/*.png
\end{lstlisting}

% =============================================================================
\section{Results}
% =============================================================================

\subsection{Exploratory Data Analysis}

\begin{figure}[H]
  \centering
  \includegraphics[width=0.85\textwidth]{01_default_rate_by_grade.png}
  \caption{Empirical default rate by Lending Club loan grade. Grade A (lowest risk)
           shows $\sim$3\% default rate; Grade G exceeds 40\%.}
  \label{fig:dr_grade}
\end{figure}

Figure~\ref{fig:dr_grade} confirms the fundamental monotonicity: default rates
increase sharply and approximately exponentially from Grade A to Grade G.
The 13$\times$ spread in default rate across grades makes \texttt{grade} the
most powerful single predictor, as confirmed by the logistic regression
coefficients in Section~\ref{sec:importance}.

\begin{figure}[H]
  \centering
  \includegraphics[width=0.85\textwidth]{02_fico_distribution.png}
  \caption{FICO score distribution by loan outcome. Defaulters skew toward lower
           FICO scores, with the distributions clearly separated.}
  \label{fig:fico}
\end{figure}

\begin{figure}[H]
  \centering
  \includegraphics[width=\textwidth]{03_key_features.png}
  \caption{DTI ratio and interest rate distributions by loan outcome. Defaulters
           show higher DTI and higher interest rates (grade-correlated).}
  \label{fig:features}
\end{figure}

\begin{figure}[H]
  \centering
  \includegraphics[width=\textwidth]{04_loan_amount_by_grade.png}
  \caption{Loan amount distribution by grade and outcome. Higher-grade borrowers
           tend to take smaller loans; the default/paid overlap is visible in all grades.}
  \label{fig:loanamt}
\end{figure}

\subsection{Classifier Performance}

\begin{figure}[H]
  \centering
  \includegraphics[width=0.87\textwidth]{05_roc_curves.png}
  \caption{ROC curves for all three classifiers. LDA and Logistic Regression
           achieve AUC $\approx 0.70$; $k$-NN underperforms due to the curse of
           dimensionality in 14-dimensional feature space.}
  \label{fig:roc}
\end{figure}

\begin{figure}[H]
  \centering
  \includegraphics[width=\textwidth]{11_model_comparison.png}
  \caption{Classifier metrics on the test set ($n=2{,}000$). All metrics
           computed at each model's chosen decision threshold.}
  \label{fig:metrics}
\end{figure}

\begin{center}
\renewcommand{\arraystretch}{1.3}
\begin{tabular}{@{}lccccc@{}}
\toprule
\textbf{Model} & \textbf{Threshold} & \textbf{Accuracy} & \textbf{Recall} & \textbf{F1} & \textbf{AUC} \\
\midrule
LDA                & 0.50 & 0.755 & 0.579 & 0.433 & 0.699 \\
$k$-NN ($k=15$)    & 0.35 & 0.690 & 0.468 & 0.246 & 0.637 \\
Logistic Reg.\ (L2) & 0.40 & 0.776 & 0.327 & 0.291 & 0.699 \\
\bottomrule
\end{tabular}
\captionof{table}{Test-set classifier performance. AUC $\approx 0.70$ is
  competitive for credit scoring; industry models typically achieve 0.70--0.85.}
\end{center}

\begin{result}
  LDA achieves the highest recall (0.579), correctly flagging more than half of
  all defaulters. Logistic Regression achieves the highest accuracy and AUC.
  $k$-NN is weakest due to the \textbf{curse of dimensionality}: in a
  14-dimensional standardized space, distances become less meaningful.
  AUC $\approx 0.70$ is consistent with published Lending Club benchmarks
  before feature engineering and ensemble methods are applied.
\end{result}

\subsection{Feature Importance}
\label{sec:importance}

\begin{figure}[H]
  \centering
  \includegraphics[width=0.87\textwidth]{06_feature_importance.png}
  \caption{Logistic regression coefficient magnitudes on standardized features.
           \texttt{grade} and \texttt{fico} dominate; delinquency history
           and interest rate follow.}
  \label{fig:importance}
\end{figure}

The top predictors in order are: (1) \textbf{grade} --- encodes the lender's
own risk assessment and strongly correlates with default; (2) \textbf{fico} ---
external credit score; (3) \textbf{delinq\_2yrs} --- payment history;
(4) \textbf{int\_rate} --- correlated with grade but adds independent signal;
(5) \textbf{pub\_rec} --- public record derogatory marks.

\subsection{Decision Threshold Analysis}

\begin{figure}[H]
  \centering
  \includegraphics[width=0.87\textwidth]{07_threshold_analysis.png}
  \caption{Precision--recall tradeoff as a function of decision threshold for
           Logistic Regression. A threshold of 0.40 (rather than 0.50) is chosen
           to increase recall at the cost of some precision.}
  \label{fig:threshold}
\end{figure}

In credit risk, \textbf{false negatives} (approving a loan that will default)
are more costly than false positives (rejecting a borrower who would have paid).
A threshold of 0.40 improves recall from 0.18 (at 0.50) to 0.33, at the cost
of a precision reduction from 0.51 to 0.45.

\subsection{Predicted PD Calibration}

\begin{figure}[H]
  \centering
  \begin{subfigure}[b]{0.49\textwidth}
    \includegraphics[width=\textwidth]{08_pd_by_grade.png}
    \caption{Model PD vs empirical default rate per grade.}
  \end{subfigure}
  \hfill
  \begin{subfigure}[b]{0.49\textwidth}
    \includegraphics[width=\textwidth]{12_calibration.png}
    \caption{Calibration plot: predicted vs observed per decile.}
  \end{subfigure}
  \caption{Logistic Regression PD calibration. The model underestimates default
           rates at high PD (Grades F/G), a known effect of L2 regularisation
           which shrinks probabilities toward the mean.}
  \label{fig:calibration}
\end{figure}

% =============================================================================
\section{Portfolio Credit Risk}
% =============================================================================

\subsection{Monte Carlo Loss Distribution}

\begin{figure}[H]
  \centering
  \includegraphics[width=0.90\textwidth]{09_loss_distribution.png}
  \caption{Simulated portfolio loss distribution (one-factor Gaussian copula,
           100,000 scenarios). The distribution is right-skewed due to systematic
           risk: bad economic states cause correlated defaults across all borrowers.}
  \label{fig:loss_dist}
\end{figure}

\begin{center}
\renewcommand{\arraystretch}{1.3}
\begin{tabular}{@{}lrl@{}}
\toprule
\textbf{Risk Metric} & \textbf{Value} & \textbf{Interpretation} \\
\midrule
Expected Loss (EL)          &  9.60\% & Average portfolio loss across all scenarios \\
Unexpected Loss (UL, std)   &  4.16\% & One standard deviation of loss \\
VaR 95\%                    & 17.28\% & Exceeded in only 5\% of scenarios \\
VaR 99\%                    & 21.29\% & Exceeded in only 1\% of scenarios \\
CVaR 95\%                   & 19.74\% & Mean loss when exceeding VaR 95\% \\
CVaR 99\%                   & 23.36\% & Mean loss when exceeding VaR 99\% \\
Basel II Capital Requirement& 16.25\% & Regulatory capital (99.9\% confidence) \\
\bottomrule
\end{tabular}
\captionof{table}{Portfolio risk metrics from Monte Carlo simulation (LGD = 45\%,
  2{,}000-loan portfolio, PD from Logistic Regression).}
\end{center}

\begin{result}
  The economic capital required to absorb unexpected losses (VaR 99\% $-$ EL
  $= 21.3\% - 9.6\% = 11.7\%$ of portfolio value) is closely approximated by
  the Basel II formula (16.3\%), which targets the 99.9\% confidence level. The
  gap between VaR and CVaR (21.3\% vs 23.4\% at 99\%) reflects the heavy right
  tail: when the systematic factor is sufficiently negative, correlated defaults
  push losses well beyond the VaR threshold.
\end{result}

\subsection{Basel II Capital Requirement}

\begin{figure}[H]
  \centering
  \includegraphics[width=\textwidth]{10_basel_capital.png}
  \caption{Left: Basel II capital requirement $K(p)$ as a function of PD (LGD=45\%).
           Right: Asset correlation $\rho(p)$ from the IRB formula. Note that
           $\rho$ decreases with PD --- riskier borrowers are assumed less
           correlated with the economic cycle.}
  \label{fig:basel}
\end{figure}

The Basel II formula (Eq.~\eqref{eq:basel}) has two key properties visible in
Figure~\ref{fig:basel}: $K(p)$ is non-monotone in $p$ (peaking around $p \approx
15\%$ and declining thereafter) and $\rho(p)$ is decreasing. This reflects the
empirical finding that high-PD borrowers (speculative grade) are more idiosyncratic,
while investment-grade borrowers default primarily in systemic crises.

% =============================================================================
\section{Conclusions \& Extensions}
% =============================================================================

\begin{center}
\renewcommand{\arraystretch}{1.4}
\begin{tabular}{@{}>{\bfseries}p{4.5cm}p{9.5cm}@{}}
\toprule
\textbf{Finding} & \textbf{Implication} \\
\midrule
LDA/LogReg achieve AUC $\approx 0.70$ &
  Comparable to published Lending Club baselines; ensemble methods
  (Random Forest, XGBoost) typically add 5--10 AUC points. \\
Grade and FICO dominate &
  Lender risk-grade is the single strongest predictor, consistent with
  the design intent of the Lending Club grading system. \\
$k$-NN underperforms &
  Curse of dimensionality in 14D standardized space; performance
  improves with lower-dimensional PCA projections. \\
CVaR $>$ VaR &
  The loss tail is heavy; regulatory models using only VaR underestimate
  true risk in stress scenarios. \\
Basel capital $>$ MC VaR 99\% &
  Basel II targets 99.9\% confidence, so this ordering is expected
  and confirms the implementation is correct. \\
\bottomrule
\end{tabular}
\end{center}

\subsection{Possible Extensions}

\begin{itemize}[leftmargin=1.8em]
  \item \textbf{PCA pre-processing}: project features to top-$k$ principal
        components before $k$-NN to mitigate dimensionality effects.
  \item \textbf{Class imbalance}: apply SMOTE or cost-sensitive learning
        to improve recall without lowering the threshold.
  \item \textbf{Survival analysis}: model \emph{time to default} rather than
        binary default, using a Cox proportional hazards model.
  \item \textbf{Stressed VaR}: shift the systematic factor $Z$ to model
        economic recession scenarios and compare with historical loss data.
  \item \textbf{Platt scaling}: recalibrate predicted probabilities post-training
        to correct the L2 shrinkage bias observed in Figure~\ref{fig:calibration}.
\end{itemize}

% ── References ────────────────────────────────────────────────────────────────
\begin{thebibliography}{9}
\bibitem{basel2}
  Basel Committee on Banking Supervision,
  \textit{International Convergence of Capital Measurement and Capital Standards
  (Basel II)}, BIS, 2006.

\bibitem{altman}
  E.~I.~Altman,
  ``Financial Ratios, Discriminant Analysis and the Prediction of Corporate Bankruptcy,''
  \textit{Journal of Finance}, 23(4):589--609, 1968.

\bibitem{thomas}
  L.~Thomas, D.~Edelman, J.~Crook,
  \textit{Credit Scoring and Its Applications}, SIAM, 2002.

\bibitem{mcneil}
  A.~McNeil, R.~Frey, P.~Embrechts,
  \textit{Quantitative Risk Management}, Princeton University Press, 2015.

\bibitem{hastie}
  T.~Hastie, R.~Tibshirani, J.~Friedman,
  \textit{The Elements of Statistical Learning}, 2nd ed.\ Springer, 2009.

\bibitem{stein}
  R.~Stein,
  ``Benchmarking Default Prediction Models,''
  \textit{Journal of Risk Model Validation}, 1(1), 2007.
\end{thebibliography}

\end{document}
